\section*{Hoofdstuk \ref{ch:g11:lenses-and-eye} --- Lenzen, beelden en het oog}

\begin{solution}{exo:g11:lenses-and-eye:1}
$C = 1/0.50 = +2.0$~D; $1/(-0.20) = -5.0$~D; $1/0.0040 = +250$~D. Een
\omterm{def:g11:lenses-and-eye:lens}{lens} van $+8.0$~D: $f' = 1/8.0 = \qty{0.125}{m} = \qty{12.5}{cm}$.
\end{solution}

\begin{solution}{exo:g11:lenses-and-eye:2}
\omterm{def:g11:lenses-and-eye:lens}{Divergerend} ($f' < 0$). Nee: voor elke voorwerpsafstand $d$ ligt
$\gamma = f'/(f' - d)$ (aangetoond in
\cref{exo:g11:lenses-and-eye:15}) tussen $0$ en $1$ --- het beeld is
altijd \omterm{def:g11:lenses-and-eye:images}{virtueel}, rechtop en \emph{verkleind}. De \omterm{def:g10:pressure:pressure}{druk} verschijnt rechtop
en gekrompen, nooit vergroot.
\end{solution}

\begin{solution}{exo:g11:lenses-and-eye:3}
Het voorwerp staat op $2f'$; de centrale straal en de evenwijdige straal
snijden elkaar \qty{20}{cm} achter de \omterm{def:g11:lenses-and-eye:lens}{lens}. Controle:
$1/\overline{OA'} = 1/10 - 1/20 = 1/20$, dus
$\overline{OA'} = +\qty{20}{cm}$ en $\gamma = 20/(-20) = -1$: \omterm{def:g11:lenses-and-eye:images}{reëel},
omgekeerd, even groot --- zoals getekend.
\end{solution}

\begin{solution}{exo:g11:lenses-and-eye:4}
$1/\overline{OA'} = 1/8.0 - 1/24 = 1/12$:
$\overline{OA'} = +\qty{12}{cm}$, $\gamma = 12/(-24) = -0.50$. Een \omterm{def:g11:lenses-and-eye:images}{reëel
beeld} \qty{12}{cm} achter de \omterm{def:g11:lenses-and-eye:lens}{lens}, omgekeerd en half zo groot.
\end{solution}

\begin{solution}{exo:g11:lenses-and-eye:5}
\omterm{def:g11:lenses-and-eye:eye}{Hoornvlies} en \omterm{def:g11:lenses-and-eye:eye}{ooglens} vormen samen de \omterm{def:g11:lenses-and-eye:lens}{convergerende lens}; het \omterm{def:g11:lenses-and-eye:eye}{netvlies} is
het scherm. De \omterm{def:g11:lenses-and-eye:eye}{accommodatie} verandert de \omterm{def:g11:lenses-and-eye:focal}{brandpuntsafstand} (de
\omterm{def:g11:lenses-and-eye:vergence}{lenssterkte}) van de \omterm{def:g11:lenses-and-eye:eye}{ooglens}; de afstand tussen \omterm{def:g11:lenses-and-eye:lens}{lens} en \omterm{def:g11:lenses-and-eye:eye}{netvlies} kan niet
veranderen. \omterm{def:g11:lenses-and-eye:eye}{Nabijpunt}: het dichtstbijzijnde punt dat bij volledige
\omterm{def:g11:lenses-and-eye:eye}{accommodatie} scherp wordt gezien --- \qty{25}{cm} voor het \omterm{def:g11:lenses-and-eye:eye}{normale oog}.
\end{solution}

\begin{solution}{exo:g11:lenses-and-eye:6}
Landschap: de sensor in het brandvlak, \qty{50}{mm} van de \omterm{def:g11:lenses-and-eye:lens}{lens}. Op
\qty{3.0}{m}: $1/\overline{OA'} = 1/50 - 1/3000 = 59/3000$, dus
$\overline{OA'} \approx \qty{50.8}{mm}$. Verplaatsing: ongeveer
\qty{0.8}{mm}.
\end{solution}

\begin{solution}{exo:g11:lenses-and-eye:7}
$1/\overline{OA'} = 1/6.0 - 1/5.0 = -1/30$:
$\overline{OA'} = \qty{-30}{cm}$ --- \omterm{def:g11:lenses-and-eye:images}{virtueel}, rechtop, met
$\gamma = (-30)/(-5.0) = +6.0$: zes keer de postzegel, \qty{30}{cm} boven
de \omterm{def:g11:lenses-and-eye:lens}{lens}.
\end{solution}

\begin{solution}{exo:g11:lenses-and-eye:8}
$1/f' = 1/40 + 1/40 = 1/20$: $f' = \qty{20}{cm}$;
$\gamma = 40/(-40) = -1$. De symmetrie dwingt
$\abs{\overline{OA'}} = \abs{\overline{OA}}$ af, en dus
$\abs{\gamma} = 1$; het beeld is \omterm{def:g11:lenses-and-eye:images}{reëel} en omgekeerd, dus $\gamma = -1$:
precies even groot.
\end{solution}

\begin{solution}{exo:g11:lenses-and-eye:9}
$1/\overline{OA'} = -1/25 - 1/50 = -3/50$:
$\overline{OA'} \approx \qty{-16.7}{cm}$,
$\gamma = (-16.7)/(-50) = +0.33$: \omterm{def:g11:lenses-and-eye:images}{virtueel}, rechtop en een derde van de
grootte. De bril van een bijziende is \omterm{def:g11:lenses-and-eye:lens}{divergerend}: alles wat je erdoor
ziet --- ook, van buitenaf, de ogen van de drager --- is verkleind.
\end{solution}

\begin{solution}{exo:g11:lenses-and-eye:10}
Het beeld van oneindig moet op het \omterm{def:g11:lenses-and-eye:eye}{vertepunt} liggen:
$f' = \qty{-0.40}{m}$, dus $C = -2.5$~D. Het beeld van de berg ligt
\qty{40}{cm} voor het glas --- precies op het \omterm{def:g11:lenses-and-eye:eye}{vertepunt}, dat het
ontspannen oog scherp ziet met nul \omterm{def:g11:lenses-and-eye:eye}{accommodatie}.
\end{solution}

\begin{solution}{exo:g11:lenses-and-eye:11}
$C = 1/\overline{OA'} - 1/\overline{OA} = 1/(-0.80) - 1/(-0.25)
= -1.25 + 4.00 = +2.75$~D.
\end{solution}

\begin{solution}{exo:g11:lenses-and-eye:12}
\emph{1.} $\gamma = \overline{OA'}/\overline{OA} = -50$ geeft
$\overline{OA'} = -50\,\overline{OA}$; met
$\overline{OA'} - \overline{OA} = \qty{5.1}{m}$ volgt
$-51\,\overline{OA} = \qty{5.1}{m}$: $\overline{OA} = \qty{-0.10}{m}$ en
$\overline{OA'} = +\qty{5.0}{m}$.

\emph{2.} $1/f' = 1/5.0 + 1/0.10 = 10.2$: $f' \approx \qty{9.8}{cm}$.

\emph{3.} $50 \times \qty{24}{mm} = \qty{1.2}{m}$, omgekeerd --- en
daarom gaan dia's ondersteboven in het toestel.
\end{solution}

\begin{solution}{exo:g11:lenses-and-eye:13}
\emph{1.} $C_\infty = 1/0.017 \approx 59$~D;
$C_{25} = 1/0.017 + 1/0.25 \approx 63$~D. \omterm{def:g10:signals-and-waves:amplitude}{Amplitude}:
$1/0.25 = 4.0$~D.

\emph{2.} \omterm{def:g10:signals-and-waves:amplitude}{Amplitude} $= 1/1.0 = 1.0$~D: een kwart ($25\%$) van de
$4.0$~D van het \omterm{def:g11:lenses-and-eye:eye}{normale oog} blijft over.
\end{solution}

\begin{solution}{exo:g11:lenses-and-eye:14}
\emph{1.} $f' = 0.25/3 \approx \qty{0.083}{m} = \qty{8.3}{cm}$.

\emph{2.} $\overline{OA'} = \qty{-25}{cm}$:
$1/\overline{OA} = -1/25 - 3/25 = -4/25$, dus
$\overline{OA} = \qty{-6.25}{cm}$ en $\gamma = (-25)/(-6.25) = +4.0$ ---
beter dan de aangeprezen $\times 3$, die het beeld op oneindig
veronderstelt; het naar het \omterm{def:g11:lenses-and-eye:eye}{nabijpunt} halen levert één \omterm{def:g10:orders-of-magnitude:unit}{eenheid} winst op.
\end{solution}

\begin{solution}{exo:g11:lenses-and-eye:15}
$1/\overline{OA'} = 1/f' - 1/d = (d - f')/(f'd)$, dus
$\overline{OA'} = f'd/(d - f')$ en
$\gamma = \overline{OA'}/(-d) = f'/(f' - d)$. Als $d > f'$: dan is
$\overline{OA'} > 0$ (\omterm{def:g11:lenses-and-eye:images}{reëel}) en $\gamma < 0$ (omgekeerd). Als
$0 < d < f'$: dan is $\overline{OA'} < 0$ (\omterm{def:g11:lenses-and-eye:images}{virtueel}) en
$\gamma = f'/(f' - d) > 1$: rechtop en vergroot --- de loepstand.
\end{solution}

\begin{solution}{pb:g11:lenses-and-eye:1}
\textbf{1.} $1/\overline{OA'} = 1/20 - 1/60 = 1/30$:
$\overline{OA'} = +\qty{30}{cm}$, $\gamma = -0.50$ --- \omterm{def:g11:lenses-and-eye:images}{reëel}, omgekeerd,
half zo groot.

\textbf{2.} $\overline{OA'} = +\qty{60}{cm}$, $\gamma = -2$: de standen
\qty{30}{cm} en \qty{60}{cm} hebben van rol gewisseld en de \omterm{def:g11:lenses-and-eye:magnification}{vergrotingen}
zijn elkaars omgekeerde --- omgekeerd licht volgt dezelfde weg.

\textbf{3.} $1/\overline{OA'} = 1/20 - 1/10 = -1/20$:
$\overline{OA'} = \qty{-20}{cm}$, $\gamma = +2$ --- \omterm{def:g11:lenses-and-eye:images}{virtueel}, rechtop,
verdubbeld.

\textbf{4.} Vanaf oneindig begint het beeld in het brandvlak; naarmate
het voorwerp $F$ nadert, wijkt het reële, omgekeerde beeld naar oneindig
en groeit het; binnen $f'$ wordt het \omterm{def:g11:lenses-and-eye:images}{virtueel}, rechtop en vergroot, vóór
de \omterm{def:g11:lenses-and-eye:lens}{lens}.

\textbf{5.} $1/\overline{OA} \approx 0$ laat
$1/\overline{OA'} = 1/f'$ over: $\overline{OA'} = f'$, het brandvlak.

\textbf{6.} In het brandvlak, \qty{50}{mm} achter de \omterm{def:g11:lenses-and-eye:lens}{lens}.

\textbf{7.} $1/\overline{OA'} = 1/50 - 1/1000 = 19/1000$:
$\overline{OA'} \approx \qty{52.6}{mm}$; verplaatsing
$\approx \qty{2.6}{mm}$.

\textbf{8.} $1/\overline{OA} = 1/55 - 1/50 = -1/550$: de kleinste
scherpstelafstand is \qty{550}{mm} = \qty{55}{cm}.

\textbf{9.} $\gamma = 55/(-550) = -0.10$: het gezicht wordt afgebeeld op
\qty{20}{mm}, net binnen de sensor van \qty{24}{mm}.

\textbf{10.} $\gamma = -1$ dwingt $\overline{OA'} = -\overline{OA}$ af,
dus $2/\overline{OA'} = 1/f'$: het voorwerp \qty{100}{mm} ($= 2f'$)
ervoor, de sensor \qty{100}{mm} erachter --- het dubbele van de
landschapsstand, ver voorbij het mechanisme van \qty{55}{mm};
macro-objectieven leveren die extra verplaatsing.

\textbf{11.} Alles wat dichterbij ligt dan haar \omterm{def:g11:lenses-and-eye:eye}{nabijpunt} (\qty{1.0}{m})
is wazig: gestrekte armen brengen de bladzijde zo dicht mogelijk bij
\qty{1.0}{m} --- scherp, zij het klein.

\textbf{12.} $C = 1/(-1.0) - 1/(-0.25) = -1.0 + 4.0 = +3.0$~D.

\textbf{13.} $\gamma = (-1.0)/(-0.25) = +4.0$: vier keer groter, maar
ook vier keer verder --- dezelfde hoekgrootte; de winst zit niet in de
grootte maar in de \emph{scherpte}, want het beeld ligt nu op haar
\omterm{def:g11:lenses-and-eye:eye}{nabijpunt}.

\textbf{14.} Scherp vraagt om een beeld voorbij \qty{1.0}{m}: van de
bladzijde op \qty{25}{cm} (beeld op \qty{1.0}{m}) tot de bladzijde op
$f' = 1/3.0 \approx \qty{0.33}{m}$ (beeld op oneindig). Een ver voorwerp
wordt \qty{33}{cm} \emph{achter} de bril afgebeeld --- een \omterm{def:g11:lenses-and-eye:images}{reëel beeld}
waar het oog niets mee kan: de verte wordt wazig.

\textbf{15.} Vandaar bifocale glazen: leessterkte in de onderste helft
van het glas, geen (of de vertecorrectie) in de bovenste helft.

\textbf{16.} $1/\overline{OA'} = 1/5.0 - 1/4.0 = -1/20$:
$\overline{OA'} = \qty{-20}{cm}$, \omterm{def:g11:lenses-and-eye:images}{virtueel}, rechtop, $\gamma = +5.0$.

\textbf{17.} $1/\overline{OA} = -1/25 - 1/5.0 = -6/25$:
$\overline{OA} \approx \qty{-4.2}{cm}$; $\gamma = (-25)/(-25/6) = +6.0$.

\textbf{18.} Steen in $F$: beeld op oneindig, bekeken met nul
\omterm{def:g11:lenses-and-eye:eye}{accommodatie} --- urenlang keuren zonder vermoeide ogen.

\textbf{19.} $C_\infty = 1/0.017 \approx 59$~D;
$C_{25} = 1/0.017 + 1/0.25 \approx 63$~D: ongeveer $4$~D \omterm{def:g11:lenses-and-eye:eye}{accommodatie}.

\textbf{20.} $f' = 1/58.8 \approx \qty{17.0}{mm}$ tegenover
$1/62.8 \approx \qty{15.9}{mm}$: ongeveer \qty{1.1}{mm}, zo'n $6\%$ ---
de ingebouwde zoom die het fototoestel met \qty{2.6}{mm} verplaatsing
nabootste en die grootmoeder met $+3$~D oplapte.
\end{solution}
